% sec-04: the operation and its governance (surgical lead).  Sources:
% onpremwhippl protocol, phase1ind, tripartisan-llm-support.md.
\section{The Operation and Its Governance}

The intervention is a staged eight-arm robotic pancreaticoduodenectomy with
perioperative daraxonrasib and an on-premises large language model whose output
is advisory only \cite{onpremwhippl,phase1ind}. Three boundaries make the study
a governed system rather than an autonomous one.

\textbf{The model has no path to the arms.} Advisory output is rendered to a
display. No generated token reaches a motion controller, and no software in the
advisory path holds write access to the robot's command channel. The operating
surgeon approves every motion.

\textbf{Stopping is bounded and measured.} Arm-level stop is specified at 3
milliseconds and system-wide stop at 500 milliseconds. Both are measured on the
bench before the first patient and re-measured at each cohort boundary.

\textbf{Every artifact is verified before use.} The model runs on premises,
pinned to a repository commit, and logs a timestamped record for each
recommendation. Conversion to open operation ordered for patient safety is
recorded as clinically appropriate; conversion caused by device, advisory, or
integration failure is counted separately against the feasibility endpoint.

\begin{apptable}
\begin{tabularx}{\textwidth}{>{\raggedright\arraybackslash}p{3.55cm} >{\raggedright\arraybackslash}p{3.15cm} >{\raggedright\arraybackslash}X}
\headrow \thead{Parameter} & \thead{Value} & \thead{Basis}\\
\midrule
Population & Resectable or borderline-resectable KRAS G12-mutated PDAC & Phase 1 protocol \\
Design and size & Open-label, single-arm, 3+3 dose escalation, up to 18 treated participants & Phase 1 protocol and IND \\
Co-primary endpoints & 30-day device- or procedure-related SAEs; DLT and MTD or RP2D; protocol tasks completed without unsafe conversion & Phase 1 protocol \\
Follow-up & 28-day screening; 30-day and 90-day safety; 24-month overall survival at one qualified site, agreement pending & Phase 1 protocol \\
\end{tabularx}
\end{apptable}
\tabcap{Trial parameters carried unchanged from the published Phase 1 protocol
and the Investigational New Drug application, so the award is evaluated against
a protocol that already exists rather than one to be written later.}

\vspace{0.1em}
\noindent\footnotesize\textbf{Positioning note.} \textit{First in human refers
to the integrated surgical and advisory workflow, not to daraxonrasib, which is
investigational and already in Phase 3 evaluation. The robotic configuration is
specified by manufacturer, model, cleared configuration, arm count, and software
version at the site agreement, and no drug supply or letter of authorization is
in place with the agent's developer.}\normalsize\par
